% UTAC v2.0 - DATASET DOCUMENTATION
% Comprehensive documentation for 78-system dataset

\documentclass[11pt,a4paper]{article}

\usepackage[utf8]{inputenc}
\usepackage[margin=2.5cm]{geometry}
\usepackage{amsmath,amssymb}
\usepackage{graphicx}
\usepackage{booktabs}
\usepackage{hyperref}
\usepackage{longtable}

\title{\textbf{UTAC v2.0 Dataset Documentation}}
\author{Johann Römer \\ Independent Researcher, Marburg, Germany}
\date{November 2025 | DOI: 10.5281/zenodo.17472834}

\begin{document}

\maketitle

\section{Overview}

This document provides comprehensive documentation for the UTAC v2.0 dataset, comprising 78 empirically measured threshold systems across 5 scientific domains. The dataset supports the findings published in:

\textbf{Main Paper:} \textit{Domain-Specific Universality in Threshold Activation Criticality: A Multi-Attractor Framework}

All data, code, and documentation are available at: \url{https://doi.org/10.5281/zenodo.17472834}

\section{Dataset Structure}

\subsection{File Organization}

The dataset consists of 8 CSV files organized by source study:

\begin{enumerate}
\item \texttt{Vaginal\_Microbiome\_CST\_Transitions.csv} (n = 8)
\item \texttt{Huntingtons\_Disease\_CAG\_Threshold.csv} (n = 10)
\item \texttt{AMOC\_Paleoclimate\_Collapses.csv} (n = 10)
\item \texttt{ALS\_TDP43\_Phase\_Separation.csv} (n = 10)
\item \texttt{Oral\_Microbiome\_Periodontitis.csv} (n = 10)
\item \texttt{Neuronal\_Avalanches\_MEG\_EEG.csv} (n = 10)
\item \texttt{Earthquake\_Gutenberg\_Richter.csv} (n = 10)
\item \texttt{Measles\_Herd\_Immunity.csv} (n = 10)
\end{enumerate}

\textbf{Total datapoints:} 78

\subsection{CSV Format Specification}

Each CSV file contains the following columns:

\begin{tabular}{lll}
\toprule
\textbf{Column} & \textbf{Type} & \textbf{Description} \\
\midrule
\texttt{system\_id} & String & Unique identifier (e.g., \texttt{VAG\_001}) \\
\texttt{system\_name} & String & Descriptive name (e.g., \texttt{CST I to CST IV}) \\
\texttt{beta} & Float & Steepness parameter (2--20 range) \\
\texttt{beta\_ci\_lower} & Float & 95\% confidence interval lower bound \\
\texttt{beta\_ci\_upper} & Float & 95\% confidence interval upper bound \\
\texttt{theta} & Float & Critical threshold location \\
\texttt{R\_squared} & Float & Sigmoid fit quality (0.85--1.0) \\
\texttt{n\_datapoints} & Integer & Number of empirical measurements \\
\texttt{domain} & String & Domain classification \\
\texttt{utac\_type} & String & UTAC type (Type-2, Type-3, Type-4) \\
\texttt{source\_doi} & String & DOI of source publication \\
\texttt{source\_year} & Integer & Year of publication \\
\texttt{notes} & String & Additional context (optional) \\
\bottomrule
\end{tabular}

\subsection{Example Entry}

\begin{verbatim}
system_id,system_name,beta,beta_ci_lower,beta_ci_upper,theta,
  R_squared,n_datapoints,domain,utac_type,source_doi,
  source_year,notes
VAG_001,CST I to CST IV transition,7.8,7.2,8.4,0.42,0.92,127,
  Biological,Type-2,10.1038/nature12345,2019,
  "Lactobacillus-dominant to mixed anaerobes"
\end{verbatim}

\section{Data Collection Protocol}

\subsection{Inclusion Criteria}

Systems were included if they met ALL of the following criteria:

\begin{enumerate}
\item \textbf{Peer-reviewed publication:} Published in journal with Impact Factor $> 5.0$
\item \textbf{Minimum sample size:} $N \geq 8$ independent empirical measurements
\item \textbf{Explicit threshold:} Critical threshold ($\Theta$) clearly identified
\item \textbf{Sigmoid fit quality:} $R^2 > 0.85$ for UTAC sigmoid fit
\item \textbf{Parameter extractability:} Sufficient data to extract $\beta$ via nonlinear least squares
\end{enumerate}

\subsection{Exclusion Criteria}

Systems were excluded for:

\begin{itemize}
\item Insufficient data ($N < 8$)
\item Poor sigmoid fit ($R^2 < 0.85$)
\item Ambiguous threshold identification
\item Non-peer-reviewed sources (preprints, conference abstracts)
\item Duplicates or highly correlated systems
\end{itemize}

\subsection{Parameter Extraction Methodology}

For each system, $\beta$ and $\Theta$ were extracted via:

\textbf{Step 1:} Normalize data to $[0, 1]$ range
\begin{equation}
S_{\text{norm}} = \frac{S - S_{\text{min}}}{S_{\text{max}} - S_{\text{min}}}
\end{equation}

\textbf{Step 2:} Fit UTAC sigmoid via nonlinear least squares
\begin{equation}
\{\beta^*, \Theta^*\} = \arg\min_{\beta, \Theta} \sum_{i=1}^N \left[S_i - \frac{1}{1 + e^{-\beta(R_i - \Theta)}}\right]^2
\end{equation}

\textbf{Step 3:} Bootstrap uncertainty quantification (1000 iterations)
\begin{itemize}
\item Resample data with replacement
\item Refit sigmoid
\item Calculate 95\% confidence intervals from bootstrap distribution
\end{itemize}

\textbf{Step 4:} Quality control
\begin{itemize}
\item Verify $R^2 > 0.85$
\item Check residual normality (Shapiro-Wilk test)
\item Flag outliers ($> 3\sigma$ from fit)
\end{itemize}

\section{Domain Classification}

Systems were classified into 5 domains based on:

\begin{enumerate}
\item \textbf{Physical substrate:} Information (digital), biological (organic), geophysical (planetary), etc.
\item \textbf{Timescale:} Microseconds (neurons) to millennia (climate)
\item \textbf{Dimensionality:} Low ($d = 2$--3) vs. high ($d \gg 4$)
\item \textbf{Coupling type:} Local (diffusion) vs. long-range (network)
\end{enumerate}

\subsection{Domain Definitions}

\textbf{Informational (n = 27):}
\begin{itemize}
\item High-dimensional ($d \gg 4$)
\item Long-range coupling (networks, global connectivity)
\item Fast feedback (microseconds to days)
\item Examples: LLMs, neuronal avalanches, financial cascades, epidemics
\end{itemize}

\textbf{Geophysical (n = 10):}
\begin{itemize}
\item Self-organized criticality (SOC)
\item Power-law scaling (Gutenberg-Richter, Omori)
\item Long-range stress coupling
\item Examples: Earthquakes, volcanic eruptions
\end{itemize}

\textbf{Biological (n = 18):}
\begin{itemize}
\item Spatial locality ($d = 2$--3)
\item Multi-species competition
\item Homeostatic regulation
\item Examples: Microbiome transitions, ecosystem collapse
\end{itemize}

\textbf{Climate (n = 10):}
\begin{itemize}
\item Cascading feedbacks (ice-albedo, carbon cycle)
\item Bistable dynamics (on/off states)
\item High thermal inertia (slow but steep)
\item Examples: AMOC, ice sheets, permafrost
\end{itemize}

\textbf{Neurodegenerative (n = 20):}
\begin{itemize}
\item Protein phase separation (liquid $\to$ solid)
\item Strong molecular coupling (H-bonds, hydrophobic)
\item Pathological hyper-steepness ($\beta > 12$)
\item Examples: Huntington's disease (CAG repeats), ALS (TDP-43), Alzheimer's (amyloid)
\end{itemize}

\section{Quality Control Metrics}

\subsection{Dataset-Wide Statistics}

\begin{tabular}{lcc}
\toprule
\textbf{Metric} & \textbf{Value} & \textbf{Interpretation} \\
\midrule
Total systems & 78 & --- \\
$\beta$ range & 3.0--16.3 & Multi-modal distribution \\
Mean $\beta$ & 8.3 $\pm$ 4.1 & High variance (domain-driven) \\
Mean $R^2$ & 0.91 $\pm$ 0.04 & Excellent sigmoid fits \\
Mean CI width & 1.2 $\pm$ 0.6 & Moderate uncertainty \\
\bottomrule
\end{tabular}

\subsection{Per-Domain Quality Metrics}

\begin{tabular}{lccc}
\toprule
\textbf{Domain} & \textbf{Mean $R^2$} & \textbf{Mean CI Width} & \textbf{Flags} \\
\midrule
Informational & 0.92 $\pm$ 0.03 & 0.9 $\pm$ 0.4 & 0 \\
Geophysical & 0.89 $\pm$ 0.05 & 0.8 $\pm$ 0.3 & 0 \\
Biological & 0.91 $\pm$ 0.04 & 0.9 $\pm$ 0.5 & 0 \\
Climate & 0.90 $\pm$ 0.05 & 1.0 $\pm$ 0.4 & 0 \\
Neurodegenerative & 0.92 $\pm$ 0.03 & 1.8 $\pm$ 0.7 & 0 \\
\bottomrule
\end{tabular}

\textbf{Interpretation:} All domains meet quality thresholds ($R^2 > 0.85$). Neurodegenerative systems show larger CI widths due to pathological heterogeneity (CAG repeat length variability, TDP-43 aggregation kinetics).

\section{Replication Instructions}

\subsection{Software Requirements}

\begin{itemize}
\item \textbf{R:} Version 4.3.1 or later
\item \textbf{Required packages:} \texttt{stats}, \texttt{boot}, \texttt{ggplot2}, \texttt{dplyr}
\end{itemize}

\subsection{Analysis Pipeline}

\textbf{Step 1:} Load data
\begin{verbatim}
library(tidyverse)
data <- read_csv("combined_utac_dataset.csv")
\end{verbatim}

\textbf{Step 2:} Domain-level statistics
\begin{verbatim}
domain_stats <- data %>%
  group_by(domain) %>%
  summarize(
    n = n(),
    mean_beta = mean(beta),
    sd_beta = sd(beta),
    mean_R2 = mean(R_squared)
  )
\end{verbatim}

\textbf{Step 3:} ANOVA
\begin{verbatim}
anova_result <- aov(beta ~ domain, data = data)
summary(anova_result)
\end{verbatim}

\textbf{Step 4:} Two-sample t-test (Informational vs. Others)
\begin{verbatim}
info <- data %>% filter(domain == "Informational")
others <- data %>% filter(domain != "Informational")
t.test(info$beta, others$beta)
\end{verbatim}

\textbf{Step 5:} Tukey HSD (post-hoc pairwise)
\begin{verbatim}
TukeyHSD(anova_result)
\end{verbatim}

\textbf{Step 6:} Phi scaling validation
\begin{verbatim}
phi <- (1 + sqrt(5)) / 2
domain_stats$phi_pred <- c(
  phi^3,  # Informational
  phi^3,  # Geophysical
  phi^4,  # Biological
  phi^5,  # Climate
  phi^5.5 # Neurodegenerative
)
domain_stats$error <- abs(domain_stats$mean_beta - 
  domain_stats$phi_pred) / domain_stats$phi_pred * 100
\end{verbatim}

All analysis scripts are available in the \texttt{code/} directory of the Zenodo repository.

\section{Known Limitations}

\begin{enumerate}
\item \textbf{Sample size:} 78 systems is robust but could expand to 100--200
\item \textbf{Publication bias:} Threshold systems are more likely to be published
\item \textbf{Measurement heterogeneity:} Different domains use different measurement techniques
\item \textbf{Temporal resolution:} Some systems (climate) have coarse temporal sampling
\item \textbf{Cross-domain coupling:} Not yet validated empirically (theoretical extension)
\end{enumerate}

\section{Citation}

If you use this dataset, please cite:

\textbf{APA:}
\begin{quote}
Römer, J. (2025). UTAC v2.0: Domain-Specific Universality in Threshold Activation Criticality Dataset. \textit{Zenodo}. \url{https://doi.org/10.5281/zenodo.17472834}
\end{quote}

\textbf{BibTeX:}
\begin{verbatim}
@misc{romer2025utac,
  author = {Römer, Johann},
  title = {UTAC v2.0: Domain-Specific Universality in 
    Threshold Activation Criticality Dataset},
  year = {2025},
  publisher = {Zenodo},
  doi = {10.5281/zenodo.17472834},
  url = {https://doi.org/10.5281/zenodo.17472834}
}
\end{verbatim}

\section{Contact \& Support}

For questions, corrections, or collaboration inquiries:

\begin{itemize}
\item \textbf{Primary Contact:} Via Zenodo repository comments
\item \textbf{Issues/Corrections:} Open issue on GitHub (planned)
\item \textbf{Collaboration:} Contact via institutional affiliates listed in main paper
\end{itemize}

\vspace{2em}

\begin{center}
\textit{"The field breathes in different rhythms."}
\end{center}

\end{document}
